% SIR compartmental flow diagram.
% Vol VI Ch 8 §8.3.
\begin{figure}[htbp]
\centering
\begin{tikzpicture}[
    box/.style={
      draw=engblack, thick, rectangle, minimum width=1.8cm, minimum height=1.2cm,
      align=center, font=\large
    },
    arrow/.style={->, thick, >=stealth, draw=engblack},
    rlabel/.style={font=\small, align=center}
  ]
  \node[box, fill=engblue!15] (S) at (0,0) {$S$};
  \node[box, fill=engamber!20] (I) at (4.5,0) {$I$};
  \node[box, fill=green!15] (R) at (9,0) {$R$};

  \draw[arrow] (S) -- (I) node[midway, above, rlabel] {$\beta S I / N$};
  \draw[arrow] (I) -- (R) node[midway, above, rlabel] {$\gamma I$};

  \node[rlabel] at (0, -1.4) {Susceptible};
  \node[rlabel] at (4.5, -1.4) {Infectious};
  \node[rlabel] at (9, -1.4) {Recovered};

  \node[rlabel, anchor=north, align=center, font=\small] at (4.5, -2.5)
    {Rates: contact-and-transmission $\beta$, recovery $\gamma$.
     Basic reproduction number $R_0 = \beta/\gamma$.};
\end{tikzpicture}
\caption[SIR compartmental flow]{The SIR compartmental model
\cite{paper:v6c08-kermack-mckendrick-1927}.
Susceptible ($S$), infectious ($I$), and recovered ($R$) compartments
are connected by mass-action transmission ($\beta SI/N$) and
constant-rate recovery ($\gamma I$). Total population $N = S + I + R$
is conserved.}
\label{fig:vol06:ch08:sir-flow}
\end{figure}
